\documentclass[11pt]{article}
\usepackage[margin=1in]{geometry}
\usepackage{booktabs}
\usepackage{graphicx}
\usepackage{amsmath}
\usepackage[hidelinks]{hyperref}
\usepackage{caption}
\usepackage{natbib}

\title{Order Permutation Test Results}
\author{}
\date{}

\begin{document}
\maketitle

\section{Overview}

We evaluate whether cognitive foundation models produce predictions that are invariant to the ordering of exchangeable context trials from Psych-101.
Under exchangeability, shuffling the order of context trials should not change the predictive distribution for a held-out target trial.
We test 16 models---8 finetuned (3 Llama-Centaur, 5 Qwentaur) and 8 base (3 Llama, 5 Qwen3)---on two experiments chosen to represent contrasting exchangeability regimes: one with independently sampled trials (primary test) and one with an adaptive sequential design (negative control).

\subsection{Experiments}

\paragraph{THINGS odd-one-out \citep{hebart2023things} --- primary test.}
Participants view a triplet of everyday objects drawn from 1{,}854 concepts and select the one that is least similar to the other two (\emph{odd-one-out}).
Triplets are sampled randomly from a combinatorial space of approximately one billion possibilities; each similarity judgment is independent of which triplets appeared earlier.
Because there is no adaptive procedure, no feedback, and no sequential dependency between trials, context trials are \textbf{exchangeable by design}: seeing ``bassinet vs.\ raspberry vs.\ detonator'' before ``dollhouse vs.\ bathtub vs.\ handcuff'' should produce the same predictions as the reverse order.
A model that has learned this task structure should be invariant to context ordering.
This experiment serves as the primary test of order invariance.

\paragraph{Intertemporal choice \citep{ruggeri2022globalizability} --- negative control.}
Participants make binary choices between an immediate monetary reward and a larger delayed reward through an \textbf{adaptive staircase procedure}.
The task begins with a choice such as ``\$500 now vs.\ \$550 in 12 months''; if the participant chooses the immediate option, the delayed reward escalates (e.g., to \$600, then \$750); when the participant switches preference, the procedure advances to the next scenario (gains to losses, magnitude changes, different time frames).
This branching titration design means each subsequent trial is determined by the previous response---trials are \textbf{not exchangeable}.
Placing a late-stage trial (e.g., ``\$5{,}000 vs.\ \$5{,}500'') before the opening ``\$500 vs.\ \$550'' trial breaks the staircase logic entirely.
A model finetuned on such data should learn to exploit trial ordering as a signal for the participant's temporal discount rate.
This experiment therefore serves as a \textbf{negative control}: we \emph{expect} context-order sensitivity, and detecting it validates that the permutation test has power to identify genuine order effects.

\subsection{Procedure}

For each (participant, target trial) pair, 50 random permutations of the context order are drawn.
The variance of each response-token probability $p_i$ across permutations is computed, averaged across tokens, then across targets per participant (768 participants for THINGS; 1{,}295 for Intertemporal choice).
Lower variance indicates stronger order invariance.

\section{Results}

\subsection{THINGS odd-one-out (primary test)}

On the primary test, finetuned models exhibit substantially lower order variance than their base counterparts (Table~\ref{tab:hebart}).
The mean variance across finetuned models is $5.25 \times 10^{-4}$, compared to $4.41 \times 10^{-3}$ for base models---an \textbf{8.4$\boldsymbol{\times}$ reduction}.
Qwentaur-0.6B achieves the lowest variance ($3.63 \times 10^{-4}$), while Qwen3-0.6B shows the highest ($7.99 \times 10^{-3}$).

All models show statistically significant non-zero variance ($p < 10^{-50}$ in every case), so no model achieves perfect order invariance.
However, the finetuned models approach near-zero variance, with all eight falling below $7 \times 10^{-4}$.
This indicates that finetuning on Psych-101 teaches the model to treat context trials as approximately exchangeable when the underlying task is genuinely unordered.

\begin{table}[ht]
\centering
\caption{Order variance on THINGS odd-one-out ($n = 768$ participants per model).}
\label{tab:hebart}
\small
\begin{tabular}{llrrr}
\toprule
Model & Type & Mean Var. & Median Var. & SE \\
\midrule
Llama-Centaur-1B  & Finetuned & .000534 & .000506 & .000016 \\
Llama-Centaur-3B  & Finetuned & .000417 & .000406 & .000005 \\
Llama-Centaur-8B  & Finetuned & .000491 & .000429 & .000028 \\
Qwentaur-0.6B     & Finetuned & .000363 & .000359 & .000004 \\
Qwentaur-1.7B     & Finetuned & .000489 & .000495 & .000005 \\
Qwentaur-4B       & Finetuned & .000682 & .000580 & .000024 \\
Qwentaur-8B       & Finetuned & .000633 & .000559 & .000015 \\
Qwentaur-14B      & Finetuned & .000593 & .000526 & .000016 \\
\midrule
Llama-3.2-1B      & Base      & .007606 & .007139 & .000092 \\
Llama-3.2-3B      & Base      & .002474 & .002256 & .000037 \\
Llama-3.1-8B      & Base      & .003003 & .002779 & .000046 \\
Qwen3-0.6B        & Base      & .007994 & .006499 & .000167 \\
Qwen3-1.7B        & Base      & .004813 & .004663 & .000052 \\
Qwen3-4B          & Base      & .002989 & .002823 & .000038 \\
Qwen3-8B          & Base      & .001826 & .001678 & .000030 \\
Qwen3-14B         & Base      & .004553 & .004176 & .000068 \\
\midrule
\multicolumn{2}{l}{\textit{Finetuned avg}} & .000525 & --- & --- \\
\multicolumn{2}{l}{\textit{Base avg}}      & .004407 & --- & --- \\
\bottomrule
\end{tabular}
\end{table}

\subsection{Intertemporal choice (negative control)}

On the negative control, the pattern reverses as expected (Table~\ref{tab:ruggeri}).
Finetuned models show \emph{higher} mean order variance ($1.76 \times 10^{-2}$) than base models ($8.69 \times 10^{-3}$)---a $2.0\times$ increase.
The most order-sensitive model is Qwentaur-8B ($3.49 \times 10^{-2}$), while Qwen3-1.7B achieves the lowest variance ($3.12 \times 10^{-3}$).

This reversal is the expected outcome for the negative control.
Finetuned models, having been trained on Psych-101, learn the sequential structure of the adaptive staircase: earlier choices in the titration reveal the participant's discount rate and constrain what later trials look like.
Shuffling this natural progression---e.g., placing an escalated high-reward trial before the initial baseline choice---disrupts the learned structure and produces larger prediction changes.
In contrast, base models lack this task-specific adaptation and produce flatter, less context-sensitive predictions that are more uniform across orderings---reflecting underfitting rather than true exchangeability.
Crucially, the fact that the permutation test detects strong order effects on this task \textbf{validates the sensitivity of the test methodology}: it confirms that genuine sequential dependencies are captured, lending credibility to the near-invariance observed for THINGS.

\begin{table}[ht]
\centering
\caption{Order variance on Intertemporal choice ($n = 1{,}295$ participants per model).}
\label{tab:ruggeri}
\small
\begin{tabular}{llrrr}
\toprule
Model & Type & Mean Var. & Median Var. & SE \\
\midrule
Llama-Centaur-1B  & Finetuned & .007585 & .003376 & .000261 \\
Llama-Centaur-3B  & Finetuned & .009052 & .003915 & .000299 \\
Llama-Centaur-8B  & Finetuned & .015548 & .009300 & .000435 \\
Qwentaur-0.6B     & Finetuned & .018901 & .013400 & .000393 \\
Qwentaur-1.7B     & Finetuned & .012772 & .006111 & .000364 \\
Qwentaur-4B       & Finetuned & .017104 & .008024 & .000527 \\
Qwentaur-8B       & Finetuned & .034860 & .034143 & .000574 \\
Qwentaur-14B      & Finetuned & .025238 & .019702 & .000557 \\
\midrule
Llama-3.2-1B      & Base      & .012424 & .012860 & .000176 \\
Llama-3.2-3B      & Base      & .005604 & .005421 & .000057 \\
Llama-3.1-8B      & Base      & .009487 & .009435 & .000101 \\
Qwen3-0.6B        & Base      & .007281 & .007175 & .000089 \\
Qwen3-1.7B        & Base      & .003123 & .003053 & .000032 \\
Qwen3-4B          & Base      & .012373 & .012018 & .000151 \\
Qwen3-8B          & Base      & .010550 & .010333 & .000137 \\
Qwen3-14B         & Base      & .008693 & .008116 & .000129 \\
\midrule
\multicolumn{2}{l}{\textit{Finetuned avg}} & .017633 & --- & --- \\
\multicolumn{2}{l}{\textit{Base avg}}      & .008692 & --- & --- \\
\bottomrule
\end{tabular}
\end{table}

\subsection{Cross-task comparison}

\begin{figure}[ht]
    \centering
    \includegraphics[width=\textwidth]{permutation_grid.pdf}
    \caption{Distribution of per-participant order variance across models and experiments.
    Each row shows one experiment: THINGS odd-one-out (top, primary test) and Intertemporal choice (bottom, negative control).
    Left panels: violin and box plots with jittered individual participants; finetuned models (left group, solid) vs.\ base models (right group, hatched).
    Right panels: empirical cumulative distribution functions; solid lines denote finetuned models, dashed lines denote base models; steeper curves closer to zero indicate greater order invariance.}
    \label{fig:grid}
\end{figure}

Figure~\ref{fig:grid} highlights the contrasting patterns across experiments.
For THINGS (top row), where triplets are sampled independently, all finetuned models cluster near zero while base models show an order of magnitude more variance---consistent with finetuning having learned that independent similarity judgments are exchangeable.
The ECDF curves confirm this: finetuned models rise sharply near zero, while base model curves spread across a wider range.

For intertemporal choice (bottom row), where the adaptive staircase creates genuine sequential dependencies, finetuned models---especially larger Qwentaur variants---display the highest variance, with Qwentaur-8B exceeding $3.4 \times 10^{-2}$.
This is the expected outcome: finetuned models have learned that the progression of choices within the titration procedure carries information about the participant's discount rate, so disrupting that progression changes predictions.
Notably, finetuned models show highly skewed distributions with prominent upper outliers, while their medians are closer to those of base models---the high mean variance is partly driven by a subset of participants whose staircase sequences are especially informative about temporal discount rate (e.g., participants who triggered multiple escalation steps before switching).

\section{Summary}

\begin{enumerate}
    \item \textbf{No model achieves perfect exchangeability.} All 16 models on both tasks yield $p \approx 0$ in one-sample $t$-tests against zero variance, confirming that context order affects predictions to some degree.

    \item \textbf{Finetuning dramatically reduces order sensitivity for THINGS (primary test).} Finetuned models show $8.4\times$ lower mean order variance than base models, indicating that training on Psych-101 teaches the model to treat independently sampled similarity judgments as approximately exchangeable.

    \item \textbf{Finetuning increases order sensitivity for Intertemporal choice (negative control).} The reversal ($2.0\times$ higher variance for finetuned models) is the expected outcome: the adaptive staircase creates genuine sequential dependencies, and finetuned models have learned to exploit trial progression as a signal for the participant's discount rate.

    \item \textbf{The permutation test is validated by the negative control.} The strong order effects detected on the adaptive staircase task confirm the test's power to identify genuine sequential dependencies, lending credibility to the near-invariance observed on the independently-structured THINGS task.

    \item \textbf{Finetuned models correctly distinguish exchangeable from non-exchangeable task structures.} Exchangeability is approximated for tasks with independently sampled trials (THINGS odd-one-out) but not for tasks where a branching titration design makes trial ordering inherently informative (intertemporal choice). Both outcomes reflect correct task-specific learning.
\end{enumerate}

\bibliographystyle{plainnat}
\begin{thebibliography}{2}
\bibitem[Hebart et~al.(2023)]{hebart2023things}
Hebart, M.~N., Contier, O., Teichmann, L., et~al. (2023).
\newblock THINGS-data, a multimodal collection of large-scale datasets for investigating object representations in human brain and behavior.
\newblock \textit{eLife}, 12, e82580.

\bibitem[Ruggeri et~al.(2022)]{ruggeri2022globalizability}
Ruggeri, K., Aky\"{u}z, S., Arami\'{c}, M., et~al. (2022).
\newblock The globalizability of temporal discounting.
\newblock \textit{Nature Human Behaviour}, 6, 1386--1397.
\end{thebibliography}

\end{document}
